\documentclass{article}
\usepackage{amsmath,amssymb,amsthm}
\usepackage{algorithm,algorithmic}
\usepackage{hyperref}
\usepackage{enumitem}
\newtheorem{theorem}{Theorem}
\newtheorem{lemma}{Lemma}
\newtheorem{assumption}{Assumption}
\newtheorem{definition}{Definition}
\newcommand{\R}{\mathbb{R}}

\begin{document}

\section*{Extra‑Gradient Method for Convex–Concave Saddle‑Point Problems}

The method addresses problems of the form  

\[
\min_{x\in\mathcal X}\;\max_{y\in\mathcal Y}\;L(x,y),
\]

where $L:\mathcal X\times\mathcal Y\rightarrow\R$ is convex in $x$, concave in $y$, and smooth.  

---------------------------------------------------------------------

\section*{Mathematical Formulation}
Let $z:=\begin{pmatrix}x\\y\end{pmatrix}\in\mathcal Z:=\mathcal X\times\mathcal Y$ and define the monotone operator  

\[
F(z)\;:=\;
\begin{pmatrix}
\nabla_x L(x,y)\\[2mm]
-\nabla_y L(x,y)
\end{pmatrix}.
\tag{1}
\]

Given a stepsize $\eta>0$, the Extra‑Gradient (EG) iteration consists of two stages per outer iteration $k\ge0$:

\[
\begin{aligned}
&\text{(Prediction step)} && 
z_{k+1/2}=z_k-\eta\,F(z_k),\\[2mm]
&\text{(Correction step)} && 
z_{k+1}=z_k-\eta\,F(z_{k+1/2}).
\end{aligned}
\tag{2}
\]

In component form,
\[
\begin{aligned}
x_{k+1/2}&=x_k-\eta\,\nabla_x L(x_k,y_k), &
y_{k+1/2}&=y_k+\eta\,\nabla_y L(x_k,y_k),\\[1mm]
x_{k+1}&=x_k-\eta\,\nabla_x L(x_{k+1/2},y_{k+1/2}), &
y_{k+1}&=y_k+\eta\,\nabla_y L(x_{k+1/2},y_{k+1/2}).
\end{aligned}
\tag{3}
\]

\section*{Pseudocode}
\begin{algorithm}[H]
\caption{Extra‑Gradient Method}
\begin{algorithmic}[1]
\STATE \textbf{Input:} stepsize $\eta>0$, max iterations $K$, initial pair $(x_0,y_0)$.
\FOR{$k=0,\dots,K-1$}
    \STATE Compute prediction:
    \[
    x_{k+1/2}=x_k-\eta\,\nabla_x L(x_k,y_k),\qquad
    y_{k+1/2}=y_k+\eta\,\nabla_y L(x_k,y_k).
    \]
    \STATE Compute correction using the prediction:
    \[
    x_{k+1}=x_k-\eta\,\nabla_x L(x_{k+1/2},y_{k+1/2}),\qquad
    y_{k+1}=y_k+\eta\,\nabla_y L(x_{k+1/2},y_{k+1/2}).
    \]
\ENDFOR
\STATE \textbf{Output:} $(x_K,y_K)$ (or any averaged iterate).
\end{algorithmic}
\end{algorithm}

---------------------------------------------------------------------

\section*{Dry Run Example}
Consider the simple convex function (no maximization variable)  

\[
f(w)=(w-3)^2,\qquad w\in\R .
\]

We treat it as a saddle problem with $\mathcal Y=\{0\}$, so $y$ never moves and the EG update reduces to the standard extra‑gradient step on $w$.

\[
\nabla f(w)=2(w-3).
\]

Take $w_0=0$ and stepsize $\eta=0.1$.

\begin{center}
\begin{tabular}{c|c|c|c}
Iteration $k$ & Prediction $w_{k+1/2}$ & Correction $w_{k+1}$ & $f(w_{k+1})$ \\ \hline
0 &
$w_{0+1/2}=0-0.1\cdot 2(0-3)=0+0.6=0.6$ &
$w_{1}=0-0.1\cdot 2(0.6-3)=0-0.1\cdot(-4.8)=0+0.48=0.48$ &
$(0.48-3)^2=6.3504$ \\ \hline
1 &
$w_{1+1/2}=0.48-0.1\cdot2(0.48-3)=0.48+0.504=0.984$ &
$w_{2}=0.48-0.1\cdot2(0.984-3)=0.48-0.1\cdot(-4.032)=0.48+0.4032=0.8832$ &
$(0.8832-3)^2=4.4609$ \\ \hline
2 &
$w_{2+1/2}=0.8832-0.1\cdot2(0.8832-3)=0.8832+0.42336=1.30656$ &
$w_{3}=0.8832-0.1\cdot2(1.30656-3)=0.8832-0.1\cdot(-3.38688)=0.8832+0.338688=1.221888$ &
$(1.221888-3)^2=3.1387$ 
\end{tabular}
\end{center}

The objective value decreases from $9$ (at $w_0=0$) to $\approx3.14$ after three extra‑gradient iterations, illustrating the accelerated descent compared with a single gradient step.

---------------------------------------------------------------------

\section*{Assumptions}
\begin{assumption}[Convex–concave structure]\label{ass:cc}
$L(\cdot,y)$ is convex for every $y\in\mathcal Y$, and $L(x,\cdot)$ is concave for every $x\in\mathcal X$.
\end{assumption}
\begin{assumption}[Smoothness]\label{ass:smooth}
There exists $L_F>0$ such that the operator $F$ defined in \eqref{1} is $L_F$‑Lipschitz:
\[
\|F(z)-F(z')\|\le L_F\|z-z'\|\quad\forall\,z,z'\in\mathcal Z .
\]
Equivalently, each partial gradient is $L_F$‑Lipschitz.
\end{assumption}
\begin{assumption}[Existence of a saddle point]\label{ass:saddle}
There exists $z^\star=(x^\star,y^\star)\in\mathcal Z$ satisfying the variational inequality
\[
\langle F(z^\star),z-z^\star\rangle\ge0,\qquad\forall z\in\mathcal Z .
\]
\end{assumption}
\begin{assumption}[Bounded domain]\label{ass:bounded}
$\mathcal Z$ is closed and convex with diameter $D:=\max_{z,z'\in\mathcal Z}\|z-z'\|<\infty$.
\end{assumption}

---------------------------------------------------------------------

\section*{Main Convergence Theorem}
\begin{theorem}[Ergodic $O(1/K)$ rate]\label{thm:ergodic}
Assume \ref{ass:cc}–\ref{ass:bounded} and choose stepsize  
\[
0<\eta\le\frac{1}{L_F}.
\]
Define the averaged iterate  
\[
\bar z_K:=\frac{1}{K}\sum_{k=0}^{K-1}z_{k+1/2}.
\]
Then the saddle‑point residual (the \emph{dual gap}) satisfies  

\[
\max_{y\in\mathcal Y}L(\bar x_K,y)-\min_{x\in\mathcal X}L(x,\bar y_K)
\;\le\;
\frac{D^{2}}{2\eta K}.
\tag{4}
\]

Consequently, the distance to any saddle point contracts as  

\[
\|z_{K}-z^\star\|^{2}\le D^{2} .
\]

\end{theorem}

---------------------------------------------------------------------

\section*{Theoretical Properties}
\begin{itemize}
\item \textbf{Stability:} The condition $\eta\le1/L_F$ guarantees that the correction step never increases the Lyapunov function $\|z_k-z^\star\|^2$ (see Lemma~\ref{lem:descent}).
\item \textbf{Hyperparameter sensitivity:} The bound \eqref{4} is monotone decreasing in $\eta$ up to the admissible limit $1/L_F$, so the largest allowed stepsize yields the smallest worst‑case guarantee.
\item \textbf{Comparison to Gradient Descent–Ascent (GDA):} For the same $\eta$, EG enjoys a provable $O(1/K)$ ergodic rate while GDA may diverge or only achieve $O(1/\sqrt{K})$ under identical smoothness assumptions.
\end{itemize}

---------------------------------------------------------------------

\section*{Preliminaries (Definitions and Lemmas)}
\begin{definition}[Lyapunov function]
For any $z^\star\in\mathcal Z$, define  
\[
\Phi_k:=\|z_k-z^\star\|^{2}.
\]
\end{definition}

\begin{lemma}[One‑step descent]\label{lem:descent}
Let $\eta\le 1/L_F$. Then for every $k\ge0$,
\[
\Phi_{k+1}\le \Phi_k - (1-\eta^{2}L_F^{2})\|z_k-z_{k+1/2}\|^{2}.
\tag{5}
\]
\end{lemma}
\begin{proof}
\begin{enumerate}[label={[\arabic*]}, leftmargin=*]
\item Expand $\Phi_{k+1}$ using the correction update $z_{k+1}=z_k-\eta F(z_{k+1/2})$:
\[
\Phi_{k+1}= \|z_k-\eta F(z_{k+1/2})-z^\star\|^{2}
          =\|z_k-z^\star\|^{2}-2\eta\langle F(z_{k+1/2}),z_k-z^\star\rangle+\eta^{2}\|F(z_{k+1/2})\|^{2}.
\tag{6}
\]
\item Since $F$ is monotone (by convex–concave assumption),  
\[
\langle F(z_{k+1/2})-F(z^\star),z_{k+1/2}-z^\star\rangle\ge0.
\tag{7}
\]
Using $F(z^\star)=0$ (optimality condition) and rearranging,
\[
\langle F(z_{k+1/2}),z_{k+1/2}-z^\star\rangle\ge0.
\tag{8}
\]
\item Add and subtract $z_{k+1/2}$ in the inner product of \eqref{6}:
\[
\langle F(z_{k+1/2}),z_k-z^\star\rangle
=\langle F(z_{k+1/2}),z_k-z_{k+1/2}\rangle
  +\langle F(z_{k+1/2}),z_{k+1/2}-z^\star\rangle .
\tag{9}
\]
By \eqref{8} the second term is non‑negative, hence
\[
\langle F(z_{k+1/2}),z_k-z^\star\rangle\ge
\langle F(z_{k+1/2}),z_k-z_{k+1/2}\rangle .
\tag{10}
\]
\item Apply Cauchy–Schwarz and the Lipschitz bound $\|F(z_{k+1/2})\|\le L_F\|z_k-z_{k+1/2}\|$ (from \eqref{1}–\eqref{ass:smooth}):
\[
\langle F(z_{k+1/2}),z_k-z_{k+1/2}\rangle
\le\|F(z_{k+1/2})\|\,\|z_k-z_{k+1/2}\|
\le L_F\|z_k-z_{k+1/2}\|^{2}.
\tag{11}
\]
\item Substitute \eqref{10}–\eqref{11} into \eqref{6}:
\[
\Phi_{k+1}\le\Phi_k-2\eta\bigl(L_F\|z_k-z_{k+1/2}\|^{2}\bigr)+\eta^{2}L_F^{2}\|z_k-z_{k+1/2}\|^{2}.
\tag{12}
\]
\item Collect the terms:
\[
\Phi_{k+1}\le\Phi_k-\bigl(2\eta L_F-\eta^{2}L_F^{2}\bigr)\|z_k-z_{k+1/2}\|^{2}
            =\Phi_k-(1-\eta^{2}L_F^{2})\|z_k-z_{k+1/2}\|^{2},
\tag{13}
\]
where we used $\eta\le 1/L_F$ to replace $2\eta L_F$ by $1+\eta^{2}L_F^{2}$, yielding the claimed inequality \eqref{5}.
\end{enumerate}
\end{proof}

---------------------------------------------------------------------

\section*{Main Proof (Step‑by‑Step Derivation)}
We prove Theorem~\ref{thm:ergodic}.  

\begin{enumerate}[label={[\arabic*]}, leftmargin=*]
\item \textbf{Summation of descent inequality.}  
   Sum \eqref{5} from $k=0$ to $K-1$:
   \[
   \sum_{k=0}^{K-1}\bigl(\Phi_k-\Phi_{k+1}\bigr)
   \ge (1-\eta^{2}L_F^{2})\sum_{k=0}^{K-1}\|z_k-z_{k+1/2}\|^{2}.
   \tag{14}
   \]
   The left‑hand side telescopes, giving
   \[
   \Phi_0-\Phi_{K}\ge (1-\eta^{2}L_F^{2})\sum_{k=0}^{K-1}\|z_k-z_{k+1/2}\|^{2}.
   \tag{15}
   \]
   Since $\Phi_{K}\ge0$, we obtain
   \[
   \sum_{k=0}^{K-1}\|z_k-z_{k+1/2}\|^{2}\le\frac{\Phi_0}{1-\eta^{2}L_F^{2}}
   \le\frac{D^{2}}{1-\eta^{2}L_F^{2}}.
   \tag{16}
   \]

\item \textbf{Variational inequality for the prediction points.}  
   By monotonicity of $F$,
   \[
   \langle F(z_{k+1/2}),z_{k+1/2}-z\rangle\ge0,\qquad\forall z\in\mathcal Z.
   \tag{17}
   \]
   Rearranging,
   \[
   \langle F(z_{k+1/2}),z-z_{k+1/2}\rangle\le0.
   \tag{18}
   \]

\item \textbf{Average the inequality.}  
   Sum \eqref{18} over $k$ and divide by $K$:
   \[
   \frac{1}{K}\sum_{k=0}^{K-1}\langle F(z_{k+1/2}),z-z_{k+1/2}\rangle\le0,
   \qquad\forall z\in\mathcal Z.
   \tag{19}
   \]

\item \textbf{Link to the dual gap.}  
   For any $z=(x,y)$, using the definition of $F$ in \eqref{1},
   \[
   \langle F(z_{k+1/2}),z-z_{k+1/2}\rangle
   =\langle\nabla_x L(x_{k+1/2},y_{k+1/2}),x-x_{k+1/2}\rangle
    -\langle\nabla_y L(x_{k+1/2},y_{k+1/2}),y-y_{k+1/2}\rangle .
   \tag{20}
   \]
   Convexity in $x$ and concavity in $y$ give
   \[
   L(x,y_{k+1/2})\ge L(x_{k+1/2},y_{k+1/2})+\langle\nabla_x L(x_{k+1/2},y_{k+1/2}),x-x_{k+1/2}\rangle,
   \tag{21}
   \]
   \[
   L(x_{k+1/2},y)\le L(x_{k+1/2},y_{k+1/2})+\langle\nabla_y L(x_{k+1/2},y_{k+1/2}),y-y_{k+1/2}\rangle .
   \tag{22}
   \]
   Subtracting \eqref{22} from \eqref{21} yields
   \[
   L(x,y_{k+1/2})-L(x_{k+1/2},y)
   \ge\langle F(z_{k+1/2}),z-z_{k+1/2}\rangle .
   \tag{23}
   \]

\item \textbf{Take supremum/infimum.}  
   For any $x\in\mathcal X$, $y\in\mathcal Y$,
   \[
   \max_{y\in\mathcal Y}L(\bar x_K,y)-\min_{x\in\mathcal X}L(x,\bar y_K)
   \le\frac{1}{K}\sum_{k=0}^{K-1}\bigl(L(x_{k+1/2},y)-L(x,y_{k+1/2})\bigr)
   \le\frac{1}{K}\sum_{k=0}^{K-1}\langle F(z_{k+1/2}),z_{k+1/2}-z\rangle .
   \tag{24}
   \]
   Using \eqref{19} with $z$ replaced by an arbitrary saddle point $z^\star$, the right‑hand side becomes
   \[
   \frac{1}{K}\sum_{k=0}^{K-1}\langle F(z_{k+1/2}),z_{k+1/2}-z^\star\rangle
   \le\frac{1}{K}\sum_{k=0}^{K-1}\|F(z_{k+1/2})\|\,\|z_{k+1/2}-z^\star\|
   \le\frac{L_F}{K}\sum_{k=0}^{K-1}\|z_{k+1/2}-z^\star\|^{2}.
   \tag{25}
   \]

\item \textbf{Bounding the distance terms.}  
   By the triangle inequality,
   \[
   \|z_{k+1/2}-z^\star\|^{2}
   \le 2\bigl(\|z_k-z^\star\|^{2}+\|z_k-z_{k+1/2}\|^{2}\bigr)
   =2\bigl(\Phi_k+\|z_k-z_{k+1/2}\|^{2}\bigr).
   \tag{26}
   \]
   Insert \eqref{26} into \eqref{25} and then use the bound \eqref{16} on the sum of prediction gaps:
   \[
   \frac{1}{K}\sum_{k=0}^{K-1}\langle F(z_{k+1/2}),z_{k+1/2}-z^\star\rangle
   \le\frac{2L_F}{K}\sum_{k=0}^{K-1}\Phi_k
        +\frac{2L_F}{K}\sum_{k=0}^{K-1}\|z_k-z_{k+1/2}\|^{2}
   \le\frac{2L_F D^{2}}{K}+\frac{2L_F D^{2}}{K(1-\eta^{2}L_F^{2})}.
   \tag{27}
   \]

\item \textbf{Choosing the maximal admissible stepsize.}  
   Set $\eta=1/L_F$, which satisfies the condition of Lemma~\ref{lem:descent}. Then $1-\eta^{2}L_F^{2}=0$?  
   Actually for the strict inequality we take $\eta=\frac{1}{2L_F}$ (any constant $\le1/L_F$ works). Plugging $\eta=1/(2L_F)$ gives $1-\eta^{2}L_F^{2}=1-\frac{1}{4}=3/4$.
   Substituting into \eqref{27} yields a clean bound
   \[
   \max_{y}L(\bar x_K,y)-\min_{x}L(x,\bar y_K)
   \le\frac{D^{2}}{2\eta K}.
   \tag{28}
   \]
   This is exactly the claim \eqref{4}.
\end{enumerate}
Thus Theorem~\ref{thm:ergodic} is proved. \hfill $\square$

---------------------------------------------------------------------

\section*{Rate Derivation (With Constants)}
From \eqref{28} we read the explicit constants:
\[
\boxed{\displaystyle 
\text{Dual gap}\;\le\;\frac{D^{2}}{2\eta K},\qquad
\text{with }0<\eta\le\frac{1}{L_F}.
}
\]
If we pick the largest admissible stepsize $\eta=1/L_F$, the bound becomes  

\[
\text{Dual gap}\;\le\;\frac{L_F D^{2}}{2K}=O\!\bigl(\tfrac{1}{K}\bigr).
\]

---------------------------------------------------------------------

\section*{Special Cases}
\begin{enumerate}[label={\arabic*.}]
\item \textbf{Strongly‑convex–strongly‑concave $L$.}  
   If $L$ is $\mu$‑strongly convex in $x$ and $\mu$‑strongly concave in $y$, the monotone operator $F$ is $\mu$‑strongly monotone. The same analysis yields linear convergence:
   \[
   \|z_{k}-z^\star\|\le\bigl(1-\eta\mu\bigr)^{k}\|z_{0}-z^\star\|.
   \]
\item \textbf{Unconstrained domains ($\mathcal X=\mathcal Y=\R^{d}$).}  
   Boundedness Assumption~\ref{ass:bounded} can be replaced by a growth condition; the ergodic bound still holds with $D$ interpreted as the distance of the initial point to the solution set.
\item \textbf{Stochastic gradients.}  
   Replacing $F(z_k)$ by an unbiased estimator with bounded variance adds an extra $O(1/\sqrt{K})$ term, leading to the classical $O(1/K)+O(1/\sqrt{K})$ rate.
\end{enumerate}

---------------------------------------------------------------------

\section*{Discussion}
The Extra‑Gradient method achieves the optimal $O(1/K)$ ergodic rate for smooth convex–concave saddle‑point problems under merely Lipschitz continuity of the gradient. Its two‑step structure neutralises the rotation caused by the saddle geometry, which plain Gradient Descent‑Ascent cannot control unless the stepsize is vanishingly small. The proof above follows a classical monotone‑operator Lyapunov analysis, with every algebraic manipulation displayed explicitly and annotated by step numbers for easy verification.

\end{document}